\documentclass[12pt]{article}
\usepackage[utf8]{inputenc}
\usepackage[T1]{fontenc}
\usepackage{hyperref}
\title{Yakan Language Resources for K-12 Curriculum Development}
\author{Compiled by Claude Sonnet 5.5 (claude-sonnet-5-5) for Andr\'as Kornai}
\date{}
\begin{document}
\maketitle

\section{Introduction}
Yakan (ISO 639-3 yka; Glottocode yaka1277) is an Austronesian language of Basilan, Philippines. Speaker counts range from about 106,000 to 130,000. Sources say it is taught in primary schools in Basilan. I did not find the specific textbooks or the current DepEd mother-tongue materials for it.

\section{Primary Educational Resources}
No primer or textbook was identified in the search results. The sources found only say that Yakan is taught in primary schools in Basilan. A search of DepEd mother-tongue materials would be the next step.

\section{Linguistic and Reference Materials}
\begin{itemize}
\item Yakan-English dictionary compiled by Dietlinde Behrens (Philippine Journal of Linguistics, 2002), online on Webonary: \url{https://www.webonary.org/yakan/?p=146819} (credits and acknowledgements). Another page of the same site: \url{https://www.webonary.work/yakan/overview/foreword/} (foreword).
\item Andrew D. Sherfan's earlier Yakan-English/English-Yakan dictionary (1972, Lamitan, Basilan), mentioned on the Webonary foreword page above.
\item The Filipinas Library catalogue has an entry that appeared in the results. I did not check its content: \url{https://www.filipinaslibrary.org.ph/biblio/53539}
\item Overview article: \url{https://en.wikipedia.org/wiki/Yakan_language}
\item Writing system overview: \url{https://omniglot.com/writing/yakan.htm}
\end{itemize}

\section{Religious and Community Resources}
According to Webonary, a Yakan New Testament was completed in 1984. I found no URL for it.

\section{Digital and Audio Resources}
\begin{itemize}
\item Online Yakan dictionary (Webonary): \url{https://www.webonary.work/yakan/?p=9444}
\end{itemize}

\section{Curriculum Development Context}
Yakan is reportedly used in primary education in Basilan, so a base for native-language instruction exists. Current curricular materials still need to be located, probably through DepEd Basilan and the Summer Institute of Linguistics Philippines. I did not verify either.
\end{document}